% hip-confidential-tier.tex
% Hanzo Improvement Proposal — The Confidential Tier.
% Number assigned when this lands in HIPs/.
\documentclass[11pt,a4paper]{article}

\usepackage[utf8]{inputenc}
\usepackage[T1]{fontenc}
\usepackage[margin=1in]{geometry}
\usepackage{amsmath,amssymb,amsthm}
\usepackage{booktabs}
\usepackage{longtable}
\usepackage{array}
\usepackage{enumitem}
\usepackage{xcolor}
\usepackage{microtype}
\usepackage[colorlinks=true,linkcolor=black,citecolor=black,urlcolor=blue]{hyperref}

\newtheorem{definition}{Definition}[section]
\newtheorem{claim}[definition]{Claim}
\newtheorem{remark}[definition]{Remark}

% Long identifiers are the norm here, so let the typewriter font break.
\newcommand{\code}[1]{\texttt{\hyphenchar\font=`\-\relax #1}}
\sloppy
\emergencystretch=3em

\newcolumntype{L}[1]{>{\raggedright\arraybackslash}p{#1}}

% A requirement identifier. Every normative sentence in this document carries
% one, so conformance can be reported against a list rather than against prose.
\newcommand{\req}[1]{\textbf{#1}}

\setlist[itemize]{itemsep=2pt,topsep=4pt}
\setlist[description]{itemsep=3pt,topsep=4pt}

\title{\textbf{The Confidential Tier}\\[6pt]
\large Evidence, verification and settlement for inference\\
on hardware Hanzo does not own\\[4pt]
\normalsize A Hanzo Improvement Proposal}

\author{Hanzo AI}
\date{5 September 2026}

\begin{document}
\maketitle

\begin{center}
\begin{tabular}{@{}ll@{}}
\toprule
Number      & assigned on merge into \code{HIPs/} \\
Type        & Standards Track \\
Category    & Core \\
Status      & Draft \\
Requires    & HIP-0027 (KMS), HIP-0095 (QoS challenges), HIP-0096 (compute\\
            & rewards), HIP-0121 (BYO fleet), HIP-0901 (Proof of AI),\\
            & HIP-1325 (\code{/v1/node}) \\
Pairs with  & \emph{Compute Settlement} \cite{hipsettle}, which specifies the ledger\\
            & this tier's receipts settle against. That proposal decides what\\
            & a claim is worth; this one decides what admits it. \\
Analysis    & \emph{Attested Inference on Operator Hardware} \cite{paper}, which this\\
            & document does not restate. Read it for why; read this for what. \\
\bottomrule
\end{tabular}
\end{center}

\begin{abstract}
\noindent
An operator who is not Hanzo runs a Hanzo image on a machine Hanzo cannot
inspect, serves inference on it, and is paid automatically. This document is the
contract that arrangement runs under: what evidence the operator produces, what a
verifier checks before believing it, what settlement pays against, and what none
of it establishes.

The companion paper \cite{paper} analyses the tier and classifies its guarantees.
It leaves three questions open that only a specification can close, and closing
them is this document's reason to exist.

First, a hardware report carries exactly one 64-byte caller field, and three
parties in the current design each believe that field is theirs: the measurement
layer shipping in \code{hanzoai/vm} binds the boot log there, the paper's flow
binds the session keys and the freshness nonce there, and the settlement path
wants the per-request receipt hash there. One field cannot carry three bindings,
and a report is generated once per boot while receipts accrue per request.
Section~\ref{sec:binding} fixes the field's contents, and routes the per-request
commitment through the receipt key named inside it rather than through the field.

Second, the assurance level a verifier assigns is currently derived, in
\code{hanzo-pqc}, from the host's own account of its hardware. A level derived
that way is an operator's claim about itself. Section~\ref{sec:levels} requires it
to be derived from the verified report, retires the second of the two overlapping
level enumerations, and states what each level does \emph{not} protect --- in
particular that a CPU environment wrapped around a passthrough accelerator
protects the control plane and not the computation.

Third, this tier is routinely described, by us as much as by anyone, in language
that belongs to a different tier. An attestation establishes that an environment
with a stated initial state was launched on a genuine part. It does not establish
that the environment ran the advertised model, and a model commitment does not
establish it either: a commitment binds a name to bytes, not bytes to a
multiplication. The claim that the advertised model produced a given output is
\emph{inherited} from assumptions about vendors, firmware and measured software,
and it fails silently when any of them fails. Section~\ref{sec:claims} makes the
distinction a conformance requirement rather than a matter of taste, with one
table that every surface reporting this tier projects from.

Section~\ref{sec:conformance} states what a conformant operator, verifier and
settlement implementation must do, and Appendix~\ref{app:tree} reports how far
the tree is from it. Nothing in this tier settles today: the on-chain evidence
backends are deployed with an empty voucher set and refuse everything, which is
the correct posture for a tier whose verifier is unfinished.
\end{abstract}

\tableofcontents

% =====================================================================
\section{Scope}
\label{sec:scope}

This document specifies the confidential tier of Hanzo's compute market: the tier
used when the served numeric path admits no exact-arithmetic check, so that
execution integrity cannot be proved and is instead inherited under stated
assumptions.

Two properties are at stake and they are not the same property.

\begin{description}[leftmargin=1.4em]
\item[Confidentiality.] The user's prompt, the returned completion, and --- where
the weights are not public --- the model, are not readable by the operator who
owns the machine they run on.
\item[Execution integrity.] The output the operator returns and bills for is the
output the advertised model produces on the given input under the advertised
decoding regime.
\end{description}

Confidential computing addresses the first directly. It addresses the second only
by implication, and the implication is conditional. This document specifies both
and labels each guarantee with the class it actually belongs to.

\paragraph{In scope.} The evidence an operator produces, the checks a verifier
performs, the release of model weights to a measured environment, the receipt
settlement pays against, and the language every surface is permitted to use about
the result.

\paragraph{Out of scope.} Exact-arithmetic verification of a forward pass, which
is HIP-0901's Proof of AI \cite{hip0901} and is preferred to this tier wherever the
numeric path allows it. Training on operator hardware, which is HIP-0901's
contribution proof. Availability and consistency, which are HIP-0095's challenges
\cite{hip0095}. Reward and bond sizing, which is HIP-0096 \cite{hip0096};
Section~\ref{sec:audit} states the inequality those numbers must satisfy and does
not set them. The identity of the node itself, which is HIP-1325 \cite{hip1325}
and its companion proposal on node identity.

\paragraph{Reading the requirements.} The key words MUST, MUST NOT, SHOULD, SHOULD
NOT and MAY are to be interpreted as in RFC 2119. Every normative sentence carries
an identifier (\req{B3}, \req{V7}, \ldots) so that an implementation can be
reported against a list. \emph{Observed} marks a statement about code that was
read, with the location given; \emph{proposed} marks a change this document asks
for. Sentences with neither marker are definitions or consequences.

% =====================================================================
\section{Terms}
\label{sec:terms}

The companion paper \cite{paper} defines measurement, attestation, trusted
computing base, model commitment, receipt and execution integrity, and this
document uses those definitions unchanged. Three further terms are needed because
the requirements below turn on them.

\begin{definition}[Report data]
The 64-byte caller-supplied field of a hardware attestation report: \code{REPORT\_DATA}
on AMD SEV-SNP \cite{snpabi} and \code{REPORTDATA} on Intel TDX \cite{tdx}. It is
supplied by software inside the environment at the moment the report is generated
and is covered by the platform's signature. Written $\mathrm{rd}$.
\end{definition}

\begin{definition}[Claim]
The 32-byte value an on-chain evidence backend judges: the thing a voucher vouches
for. Observed: \code{luxfi/standard} at \code{contracts/ai/compute/} names this
parameter \code{reportData} in both \code{CCEvidence} and \code{TEEEvidence}.
\end{definition}

\begin{remark}[Two names, two objects]
\label{rem:collision}
Report data and claim are different objects of different widths at different
layers, and the contracts currently give the second the first's name. The
collision is not cosmetic: a reader of \code{TEEEvidence} reasonably concludes
that the value being vouched for is the value inside the hardware report, and
\req{B4} below says it is not and cannot be. \emph{Proposed:} the contract
parameter is renamed \code{claim}. Until it is, a reader must hold the two apart
by position, and this document uses \emph{claim} for the chain-side value
throughout.
\end{remark}

\begin{definition}[Assurance level]
An element of the ordered set of Section~\ref{sec:levels}, describing what an
environment's evidence establishes about where a computation was protected. A
level is a conclusion a verifier reaches, never a property a host reports.
\end{definition}

% =====================================================================
\section{The binding}
\label{sec:binding}

This section is the document's substance. Everything else follows from it or
checks it.

\subsection{One field, three claimants}

A hardware report carries one report-data field. Report generation is a firmware
operation, rate-limited on both platforms, and happens once per boot; requests
arrive by the thousand against that one boot. Three bindings are currently
proposed for the same 64 bytes.

\begin{center}
\small
\begin{tabular}{@{}L{0.26\textwidth}L{0.40\textwidth}L{0.26\textwidth}@{}}
\toprule
Claimant & Binds & Scope \\
\midrule
\code{vm-measure} (observed: \code{hanzoai/vm}, \code{crates/vm-measure/src/lib.rs}) &
$\mathsf{SHA\text{-}512}$ over the launch and workload registers &
Per boot \\
The paper's flow \cite{paper} &
$\mathsf{keccak}(\text{tag} \,\|\, pk_e \,\|\, pk_r \,\|\, \mathrm{sid} \,\|\, n_v)$ &
Per boot \\
The settlement path \cite{paper} &
The receipt hash $\rho$ &
Per request \\
\bottomrule
\end{tabular}
\end{center}

The third is impossible as stated: a per-request value cannot occupy a per-boot
field. The first two are each incomplete. A binding over the boot log alone omits
the keys the session runs on and omits any freshness term, so the report is
replayable by anyone holding a copy and says nothing about which key material the
environment holds. A binding over the keys alone omits the workload log, and on
SEV-SNP the workload log has no platform counterpart --- the platform measures the
launch and nothing after it, so a launch-only report cannot speak for what the
started system was told to run.

\subsection{What the field carries}

\begin{description}[leftmargin=1.6em]
\item[\req{B1}] The report-data field MUST be
\[
\mathrm{rd} \;=\; \mathsf{SHA\text{-}512}\big(\text{``hanzo/rd/v1''} \,\|\, L \,\|\, W \,\|\, pk_e \,\|\, pk_r \,\|\, \mathrm{sid} \,\|\, n_v\big),
\]
where $L$ and $W$ are the launch and workload registers of the measurement
(48 bytes each, SHA-384, as \code{vm-measure} computes them), $pk_e$ is the
environment's X-Wing public key (1216 bytes: the X25519 share followed by the
ML-KEM-768 share), $pk_r$ its ed25519 receipt verifying key (32 bytes),
$\mathrm{sid}$ the session identifier (16 bytes) and $n_v$ the verifier's nonce
(32 bytes). $W$ is the workload register at the instant the report is requested,
which \req{B2} fixes as the end of the chain.

\item[\req{B1a}] Every term of that preimage MUST be of the fixed width stated
above, including the tag. Nothing is length-prefixed and nothing may be: a
concatenation of fixed-width fields is injective, so two distinct tuples cannot
produce one preimage, and adding length prefixes to fields that already have
lengths is a second mechanism for a property the first already has. The
consequence is a constraint on the session identifier rather than on the hash: an
implementation that wants variable-length session names carries the name
elsewhere and puts its 16-byte digest here.

\item[\req{B2}] The report MUST be requested only after the measurement chain is
complete. Nothing may be extended into either register afterwards. \emph{Observed:}
\code{Measurement::observe} already enforces this ordering by taking the binding as
its argument at the point the report is requested.

\item[\req{B3}] $sk_e$ and $sk_r$ MUST be generated inside the environment, MUST
NOT be persisted, and MUST NOT outlive the environment. A key that survives a
reboot is a key whose binding to a report is no longer implied by the report.

\item[\req{B4}] The per-request receipt MUST NOT be carried in report data. It is
bound to the hardware through $pk_r$: the report attests $pk_r$, and $sk_r$ signs
each receipt. A verifier follows the chain report $\rightarrow pk_r \rightarrow$
receipt, and \req{V12} states the link it must check.

\item[\req{B5}] Where the tier requires accelerator evidence, each device report's
nonce MUST be
$\mathsf{SHA\text{-}512}(\text{``hanzo/rd/gpu/v1''} \,\|\, \mathrm{rd} \,\|\, \text{device index})$,
truncated to the width the device's report accepts. Deriving the device nonce from
$\mathrm{rd}$ is what makes the CPU and device reports one statement; without it an
operator may present a CPU report from one machine and a device report from
another.
\end{description}

\subsection{Which digest, and why it is not two ways of doing one thing}

\begin{description}[leftmargin=1.6em]
\item[\req{B6}] Values written into the 64-byte report-data field MUST be computed
with SHA-512. Values written into a 32-byte on-chain commitment --- the claim of
Section~\ref{sec:terms}, the model commitment, HIP-0901's transcript fields --- MUST
be computed with keccak-256.
\end{description}

The two are one rule, not two: \emph{the digest is chosen by the width of the field
it fills}. A 64-byte field filled with a 32-byte digest is a field half of which is
padding, and padding inside a signed structure is space whose contents nobody
checks. A 32-byte on-chain word filled with a truncated 64-byte digest is a
gratuitous divergence from the hash the EVM and HIP-0901 already use everywhere
else. \emph{Observed:} \code{vm-measure} already reasons this way in the tree,
choosing SHA-512 for the binding for exactly this reason, and SHA-384 for the
registers because that is the digest both a TDX runtime measurement register and
the SEV-SNP launch measurement use, so nothing is re-hashed at the boundary to
hardware.

\begin{description}[leftmargin=1.6em]
\item[\req{B7}] The measurement registers MUST be SHA-384 and MUST be extended with
the rule
$R_{i+1} = \mathsf{SHA\text{-}384}(R_i \,\|\, e_i)$, $e_i = \mathsf{SHA\text{-}384}(\text{name}_i \,\|\, \texttt{0x00} \,\|\, \mathsf{SHA\text{-}384}(\text{content}_i))$,
from a register of 48 zero bytes. \emph{Observed:} implemented in
\code{vm-measure}; this requirement pins it so a second implementation cannot
diverge.

\item[\req{B8}] A measurement event MUST hash content only. The path a file was
read from MAY be recorded for a reader and MUST NOT enter the digest, so that the
same assets measure identically on every host that holds them. A measurement that
is a fingerprint of one host's filesystem cannot be predicted by a verifier, and a
measurement a verifier cannot predict is not an allowlist entry.
\end{description}

\subsection{What the chain-side nonce is not}
\label{sec:twononces}

A companion proposal on compute settlement \cite{hipsettle} binds an operator's
epoch claim to the 8-byte nonce of the 48-byte header the deployed verifier
validates, so that a quote produced for one operator's epoch cannot be presented
for another's. That binding is sound for what it does and this document does not
displace it. The two operate at different layers and neither substitutes for the
other, which is worth stating because both are called a nonce.

$n_v$ under \req{B1} is inside a \emph{hardware} report and is therefore
unforgeable by anyone without the platform: it says an environment holding these
keys exists now. The header nonce is inside a structure the verifier checks only
for a certificate chain, so it establishes that a vendor-certified key signed 48
bytes. It prevents replay of another operator's header; it does not prevent an
operator holding any vendor-certified key from minting a fresh header for a
machine running no environment at all. Anti-replay and existence are different
properties, and only the first is available from the header.

\subsection{The change this asks of the tree}

\emph{Observed:} \code{Measurement::bind} currently computes
$\mathsf{SHA\text{-}512}(L \,\|\, W)$ --- the first two terms of \req{B1} and no
others. \emph{Proposed:} it takes the session terms as arguments and computes
\req{B1}. The registers, the extension rule and the choice of digest are unchanged;
what changes is that the binding stops being a property of the boot alone and
becomes a property of the boot \emph{and} the session it will serve. Reports
produced under the current binding are not replayable-by-accident so much as
replayable-by-construction, and no verifier check can repair that after the fact.

% =====================================================================
\section{The evidence bundle}
\label{sec:bundle}

\begin{description}[leftmargin=1.6em]
\item[\req{E1}] An environment MUST publish its evidence as one document with the
fields of Table~\ref{tab:bundle}, JSON-encoded, with byte strings hex-encoded
lowercase and unpadded. \emph{Observed:} \code{vm-measure}'s measurement already
serializes this way and refuses to deserialize a document whose registers do not
fold from their events, or whose binding does not follow from its registers; a
bundle inherits both refusals.

\item[\req{E2}] A bundle MUST carry the raw platform report exactly as the platform
produced it. It MUST NOT carry a parsed, re-encoded or summarised report: parsing
belongs to the verifier, and a producer that parses on the verifier's behalf has
placed itself inside the verifier's trust base.

\item[\req{E3}] A producer MUST NOT synthesise a report. Where the platform is
absent or will not answer, the bundle MUST say so and MUST NOT name a platform.
\emph{Observed:} \code{vm-measure}'s \code{attest::status} is correct here --- a
device that will not answer yields \code{none} rather than an asserted platform.
\emph{Observed, and the reason this is a requirement:}
\code{hanzoai/node}'s \code{hanzo-bin/hanzo-node/src/security/tee\_attestation.rs}
carries simulation and development modes that return a fabricated SEV-SNP report
when the hardware is absent, and its production path falls back to simulation with
a warning rather than failing. That path MUST be removed rather than
configured off. A fabricated report that reaches a verifier which has not yet
implemented chain validation is indistinguishable from a real one.
\end{description}

\begin{table}[ht]
\centering
\small
\begin{tabular}{@{}lL{0.56\textwidth}l@{}}
\toprule
Field & Contents & Required \\
\midrule
\code{platform} & \code{sev-snp}, \code{tdx}, or \code{none} & always \\
\code{report} & The raw platform report & unless \code{none} \\
\code{chain} & The certificate chain, leaf first & unless \code{none} \\
\code{launch} & The launch log: register value and its ordered events & always \\
\code{workload} & The workload log, same shape & always \\
\code{bind} & $\mathrm{rd}$, per \req{B1} & always \\
\code{ke} & $pk_e$, the X-Wing public key & always \\
\code{kr} & $pk_r$, the receipt verifying key & always \\
\code{sid} & The session identifier & always \\
\code{nonce} & $n_v$, as the verifier issued it & always \\
\code{devices} & One device report and chain per accelerator, indexed & level $\geq$ 3 \\
\bottomrule
\end{tabular}
\caption{The evidence bundle. A field a producer cannot fill is omitted, never
invented.}
\label{tab:bundle}
\end{table}

% =====================================================================
\section{Verification}
\label{sec:verify}

\begin{description}[leftmargin=1.6em]
\item[\req{V0}] A verifier MUST perform every check below that applies to the
bundle's platform and level, and MUST refuse the bundle if any of them fails or
cannot be performed. There is no partial acceptance and no check whose failure
downgrades rather than refuses. The order is immaterial; the totality is not.
\end{description}

Each check is silent when skipped. This is the property that makes the list
normative rather than advisory: a verifier missing one does not error, it accepts.

\subsection{Every platform}

\begin{description}[leftmargin=1.6em]
\item[\req{V1}] Refuse a bundle any of whose \req{B1} terms is not the width
\req{B1a} fixes, before hashing anything. Then recompute $\mathrm{rd}$ from the
bundle's own \code{launch}, \code{workload}, \code{ke}, \code{kr}, \code{sid} and
\code{nonce}, and require it to equal the report's report-data field. A report whose
caller field says something else is a report about a different environment.

\item[\req{V2}] Require \code{nonce} to be a value this verifier generated, for
this session, and not previously accepted. A hardware report carries no
trustworthy notion of the present: SEV-SNP reports carry no timestamp, and any
timestamp a payload asserts is asserted by the party being checked. Freshness comes
from the verifier's own nonce inside the binding and from nothing else. A verifier
that does not generate $n_v$ itself, or does not check it, has accepted a report
that may be arbitrarily old --- and reports are public evidence, so an old one is
not hard to obtain.

\item[\req{V3}] Require the launch register to appear in the published measurement
log of Section~\ref{sec:log}, and require the workload register to appear in the
log entry for the workload the operator is serving.

\item[\req{V4}] Record the platform's chip identifier and report identifier against
the operator's registered identity, and refuse a bundle whose chip identifier is
already bound to a different operator. One machine presenting itself as several
independent operators is the Sybil case HIP-0096's reward economics assumes away;
the attestation is what makes that assumption checkable, and only if the
identifiers are recorded and compared.

\item[\req{V5}] Refuse a bundle whose \code{platform} is \code{none}. Such a bundle
is host-asserted evidence, which is a different and weaker object: it establishes
that a node made a claim and that the claim has not been altered, and it does not
establish that the node fed those files to the monitor, because nothing outside the
node observed the boot \cite{hanzod}.

This refuses host-asserted evidence \emph{at this tier}, and does not refuse it from
the system. The settlement proposal \cite{hipsettle} admits it at an emission
multiplier below one, which prices the difference rather than denying it, and that is
a coherent lower tier. What \req{V5} forbids is counting it as this one: an operator
whose evidence is its own signature must never be paid, reported, or displayed at a
level that implies a vendor stood behind it.
\end{description}

\subsection{AMD SEV-SNP}

\begin{description}[leftmargin=1.6em]
\item[\req{V6}] Validate the chain from the VCEK through the ASK to the AMD root,
for the chip identifier \emph{and the reported TCB version carried in the report
itself}. The VCEK is derived from both, so a report whose \code{REPORTED\_TCB} does
not match the key that signed it MUST fail. Distinguish VCEK-signed from
VLEK-signed reports and accept only what policy admits. Check revocation.

\item[\req{V7}] Require \code{MEASUREMENT} to equal the expected launch digest for
the image named in the measurement log.

\item[\req{V8}] Check the policy word: \code{DEBUG} MUST be clear, or the host may
attach a debugger to the guest and every guarantee in this document is void;
\code{MIGRATE\_MA} MUST be clear unless a migration agent is intended;
\code{SMT} MUST match deployment policy; the minimum ABI version MUST be met. Where
an ID block is used, check \code{ID\_KEY\_DIGEST} and \code{AUTHOR\_KEY\_DIGEST}.

\item[\req{V9}] Require \code{REPORTED\_TCB} and \code{GUEST\_SVN} to meet a
published floor, so that an operator cannot roll firmware back to a version with a
known break. Without a floor, every other check here still passes on a knowingly
broken platform.
\end{description}

\subsection{Intel TDX}

\begin{description}[leftmargin=1.6em]
\item[\req{V10}] Verify the quote against Intel's provisioning collateral: the
quoting enclave identity, the TCB info, and the revocation lists. Check
\code{MRTD} and \code{RTMR0..3} against the expected values for the image. Check
the TD attributes --- \code{DEBUG} clear, \code{XFAM} within policy.

\item[\req{V11}] Treat the TCB status strictly. \emph{Configuration needed} and
\emph{software hardening needed} are not \emph{up to date}. A verifier MAY accept a
degraded status only under a policy decision recorded against the bundle, and MUST
NOT accept it by default.
\end{description}

\subsection{The receipt chain}

\begin{description}[leftmargin=1.6em]
\item[\req{V12}] For every receipt, require the signature to verify under the
\emph{same} $pk_r$ that appears in an accepted bundle for the same $\mathrm{sid}$,
and require the receipt's session counter to be strictly increasing within that
session. A receipt under a key no accepted bundle names is a receipt from nowhere.

\item[\req{V13}] Require each accepted bundle to have a bounded lifetime, after
which the environment re-attests with a fresh nonce and fresh keys. Receipts under
an expired bundle MUST NOT settle. This bounds the window a compromised $sk_r$
covers; it does not detect the compromise, and Section~\ref{sec:residual} says why
nothing here does.
\end{description}

\subsection{Accelerators, at level 3 and above}

\begin{description}[leftmargin=1.6em]
\item[\req{V14}] Verify each device report against the vendor's device-identity
chain, and the reported firmware and VBIOS measurements against the published
reference integrity manifests for that device and driver. Require confidential mode
to be genuinely enabled rather than in a development or devtools setting. Require
the device nonce to equal the value \req{B5} derives. Where a partitioned device is
used, require the partition configuration to appear in the attested state rather
than being reported by the host.

\item[\req{V15}] A verifier MUST NOT assign a level above 2 on the basis of a
device the host merely reports. Only a device that produced a report verified under
\req{V14} counts toward a level.
\end{description}

% =====================================================================
\section{Measurements a verifier can interpret}
\label{sec:log}

An allowlist of measurements is a list of numbers. Whether it is also a security
control depends entirely on whether a reader outside Hanzo can map an entry back to
source.

\begin{description}[leftmargin=1.6em]
\item[\req{M1}] The image build MUST be reproducible: an independent party
rebuilding from the named source revision obtains the same bytes.

\item[\req{M2}] The mapping from image digest to expected launch measurement MUST
be computable offline from published inputs, so that a verifier derives the
expected value rather than being told it.

\item[\req{M3}] The set of $(\text{revision}, \text{digest}, \text{measurement})$
triples MUST be published in an append-only log with the lone-node promotion rule of
RFC 6962 \cite{rfc6962}, so that a measurement quietly added is visible and a
subtree cannot be presented as the whole log.

\item[\req{M4}] A measurement MUST NOT be admitted to the allowlist by an
administrative action alone. Without \req{M1}--\req{M3}, approving a measurement is
an assertion by whoever holds a key that a 32-byte value is trustworthy, and the
quorum of distinct measurements \code{TEEEvidence} enforces becomes a quorum over
values nobody outside can interpret --- which costs an attacker one approval and
buys a verifier nothing.
\end{description}

% =====================================================================
\section{Weight release}
\label{sec:release}

\begin{description}[leftmargin=1.6em]
\item[\req{K1}] Weights that are not public MUST reach the environment sealed to
$pk_e$ under X-Wing HPKE \cite{hpke}, released by the KMS \cite{hip0027} only when
every check of Section~\ref{sec:verify} has passed for the bundle naming that
$pk_e$. The same envelope HIP-0901 specifies for prompt sealing is used; there is
no second envelope.

\item[\req{K2}] A release MUST be scoped to the model set the operator is licensed
to serve, MUST carry a lifetime not exceeding the bundle's, and MUST be recorded in
the KMS audit trail as an access by a measured environment rather than by a person.

\item[\req{K3}] The environment MUST NOT write released weights back in plaintext,
to disk or to any surface the host can read.

\item[\req{K4}] A weight set whose disclosure would be unacceptable MUST NOT be
released to third-party hardware at any level. The level ladder measures cost to an
attacker who owns the machine, and that cost is finite; where the loss from
disclosure exceeds it, the control is a placement decision and not a cryptographic
one.
\end{description}

% =====================================================================
\section{The receipt and what settles}
\label{sec:receipt}

The receipt exists to make model substitution a detectable event rather than an
invisible one. It reuses HIP-0901's transcript fields so the system has one
commitment format, and adds the three bindings that format does not carry.

\begin{description}[leftmargin=1.6em]
\item[\req{R1}] For each request the environment MUST sign, under $sk_r$, a receipt
carrying the fields of Table~\ref{tab:receipt}, over the preimage
$\rho = \mathsf{keccak}(\text{``hanzo/receipt/v1''} \,\|\, R)$.

\item[\req{R2}] The sampling seed MUST be derived, not chosen:
$\mathsf{keccak}(\mathrm{sid} \,\|\, i \,\|\, n_v \,\|\, n_c)$, where $n_c$ is a
nonce the client places inside the sealed request. Every term is fixed before
generation and one of them belongs to the client, so an environment can neither
search seeds for a favourable output nor reuse a session's randomness across
requests.

\item[\req{R3}] The engine version $\kappa$ MUST record the reduction order and
batch-invariance settings in force, because an audit that replays a request cannot
reproduce it otherwise. Reduction order in a fused kernel typically depends on batch
shape, so the same model on the same input yields different logits under different
concurrent load; a receipt that does not say which regime was in force is not
replayable even by an honest party.
\end{description}

\begin{table}[ht]
\centering
\small
\begin{tabular}{@{}lL{0.62\textwidth}@{}}
\toprule
Field & Meaning \\
\midrule
$D_w$ & Model commitment: keccak Merkle root over the quantized weight bytes as
loaded, with the quantization discriminant per leaf. HIP-0901's \code{modelSpecHash}. \\
$Q$ & Quantization specification. \\
$D_t$ & Tokenizer digest. \emph{Not in HIP-0901.} \\
$\kappa$ & Engine and kernel version, per \req{R3}. \\
$s$ & Sampler specification, including the derived seed of \req{R2}. \emph{Seed not
in HIP-0901.} \\
$x_{\text{raw}}$ & Digest of the input as the user sent it, before tokenization.
\emph{Not in HIP-0901.} \\
$x_{\text{tok}}$ & Digest of the canonical input token identifiers. HIP-0901's
\code{inputCommitment}. \\
$y_{\text{tok}}, y_{\text{raw}}$ & Digests of the output token identifiers and of the
detokenized bytes. \\
$n_{\text{in}}, n_{\text{out}}$ & Token counts, which is what billing multiplies. \\
$\mathrm{sid}, i$ & Session identifier and a counter monotone within the session. \\
\bottomrule
\end{tabular}
\caption{Receipt fields. Three are new relative to HIP-0901's transcript and each
closes a substitution the transcript alone admits.}
\label{tab:receipt}
\end{table}

\paragraph{Why the tokenizer digest and the raw input.} HIP-0901 commits the input
as canonical token identifiers. A tokenizer swap changes the map from user text to
identifiers, and therefore the output, while leaving the model commitment, the
quantization specification and the kernel version untouched. The committed
identifiers do change --- but they are the identifiers \emph{the environment chose},
so a user checking a commitment against their own text has no way to know which
tokenizer to check with. Binding $D_t$ and $x_{\text{raw}}$ closes it: anyone
holding the advertised tokenizer recomputes $x_{\text{tok}}$ from $x_{\text{raw}}$
and rejects a mismatch.

\subsection{Settlement}

\begin{description}[leftmargin=1.6em]
\item[\req{S1}] A billed request MUST settle against a receipt. A request whose
receipt is absent, unsigned by a $pk_r$ named in an unexpired accepted bundle, or
bound to a model commitment outside the operator's licensed set MUST settle at zero.

\item[\req{S2}] A receipt naming a smaller model than the one advertised MUST settle
at the smaller model's rate rather than being refused, so that substitution is
priced rather than merely disputed.

\item[\req{S3}] An implementation MUST NOT fall back to a claimed usage record when
a receipt is missing. This is the requirement most likely to be undone for
operational reasons, and undoing it removes the effect of every other requirement in
this section: an operator who routes a request to a cheaper endpoint that never
enters the environment produces no receipt at all, and a usage-record fallback pays
for it anyway.

\item[\req{S4}] The claim a voucher records on chain MUST be $\rho$, and the
no-double-mint ledger MUST be keyed on the computation commitment, so one unit of
work settles once network-wide. \emph{Observed:} HIP-0901's
\code{AIComputeRegistry} already keys its ledger this way.

\item[\req{S5}] A voucher MUST NOT record a vouch until it has performed
Section~\ref{sec:verify} in full for the bundle the receipt chains to.
\emph{Observed:} \code{TEEEvidence}'s documented contract with its voucher already
requires a report-data equality check that the currently deployed verifier cannot
perform, because the payload that verifier validates is a 48-byte Lux-defined header
--- device identifier, timestamp, nonce --- with no measurement field, no policy
word, no TCB version and no report-data field. A successful verification there
establishes that a key certified under an embedded vendor root signed those 48
bytes, and no more. \emph{Proposed:} the verifier parses genuine SEV-SNP and TDX
reports, which supplies the report-data field and the measurement, policy and TCB
fields \req{V6}--\req{V11} need, none of which the header carries either.
\end{description}

% =====================================================================
\section{Assurance levels}
\label{sec:levels}

\subsection{One enumeration}

\emph{Observed:} the tree carries two enumerations of one concept.
\code{hanzo-pqc}'s \code{PrivacyTier} has five levels ordered by what is protected;
\code{CCTier} has four ordered inversely and named for mechanism.

\begin{description}[leftmargin=1.6em]
\item[\req{L1}] \code{PrivacyTier} survives. \code{CCTier} MUST be removed. Two
enumerations of one concept are two answers to one question, and the survivor is the
one whose order is monotone in protection, because that is the order a comparison
like ``at least level 3'' is written against.
\end{description}

\begin{center}
\small
\begin{tabular}{@{}cL{0.34\textwidth}L{0.40\textwidth}@{}}
\toprule
Level & Requires & Does not protect \\
\midrule
0 & Nothing & Anything \\
1 & Encryption at rest under a key the host does not hold & Anything in use \\
2 & A CPU environment attested per \req{V6}--\req{V11} & \emph{Any computation on a
discrete accelerator} \\
3 & Level 2, and every serving device attested per \req{V14} in confidential mode &
Metadata; physical attack \\
4 & Level 3 with the device inside the trust domain (TEE-I/O), removing bounce
buffers & Metadata; physical attack \\
\bottomrule
\end{tabular}
\end{center}

\begin{description}[leftmargin=1.6em]
\item[\req{L2}] A verifier MUST derive the level from verified evidence. A host's
own account of its hardware MAY determine what the host attempts and MUST NOT
determine what a verifier concludes. \emph{Observed:} \code{VendorCapabilities::max\_tier}
in \code{hanzo-pqc} derives a level from self-reported vendor features. As a
statement of what a host will try, that is correct and useful. As an input to a
verifier's conclusion it is an operator grading its own work, and no caller may use
it that way.

\item[\req{L3}] A surface reporting level 2 for a workload served on a discrete
accelerator MUST NOT describe the prompt as confidential during inference. A CPU
environment protects memory the CPU owns; a device that is not itself a trusted
environment cannot perform DMA into encrypted guest memory, so weights on load and
activations and KV cache in flight cross through shared pages and are readable by
the host. \emph{Level 2 protects the control plane and not the computation.} It is a
real step up for CPU-side work and it is not what most readers will assume it is.
\end{description}

% =====================================================================
\section{What may be claimed}
\label{sec:claims}

Every guarantee this tier produces belongs to one of four classes, and the class is
not a matter of emphasis.

\begin{description}[leftmargin=1.4em]
\item[C --- cryptographic] Holds by a stated hardness assumption or
information-theoretically, with no trust in a manufacturer.
\item[H --- hardware trust] Holds if named vendors, firmware and measured software
behave as specified.
\item[E --- economic] Holds if a rational adversary finds cheating unprofitable.
\item[O --- open] Not established by anything in this document.
\end{description}

\begin{description}[leftmargin=1.6em]
\item[\req{C1}] Any surface that reports this tier --- an API field, a dashboard, a
settlement record, a document --- MUST carry the class of the property it reports,
projected from Table~\ref{tab:class}. One table, one projection; a surface that
restates the classes in its own words will drift from them.

\item[\req{C2}] The words \emph{verified}, \emph{proven} and \emph{guaranteed} MUST
NOT be applied to a class-H property. The available honest form names the assumption:
``attested under AMD's root at TCB $\geq v$'' is a class-H statement written
correctly; ``verified inference'' is the same statement written as though it were
class C.

\item[\req{C3}] The phrase \emph{verifiable inference} MUST NOT be used for this
tier at all. Where the numeric path admits Freivalds, that path is class C and the
phrase belongs to it; where it does not, the phrase names something the system does
not have. This applies to our own documents, of which several currently violate it.

\item[\req{C4}] A surface reporting this tier's settlement MUST NOT describe the
path as post-quantum. Hardware attestation signatures are classical elliptic-curve
or RSA. \emph{Observed:} \code{TEEEvidence} already records this as the sole
classical assumption in an otherwise post-quantum tier, and records that it is
governance-removable. Until post-quantum device-identity roots exist in silicon, the
attestation leg is not post-quantum and saying otherwise misstates the strongest
claim in the system.
\end{description}

\begin{table}[ht]
\centering
\small
\begin{tabular}{@{}L{0.30\textwidth}L{0.28\textwidth}cL{0.28\textwidth}@{}}
\toprule
Property & Mechanism & Class & What breaks it \\
\midrule
Prompt confidential in transit & X-Wing HPKE to $pk_e$ & C & ML-KEM-768 \emph{and}
X25519 both broken \\
Prompt confidential during CPU execution & CPU environment isolation & H & Vendor
key, firmware, side channel, physical attack \\
Prompt confidential during device execution & Device confidential mode, level $\geq$ 3
& H & As above plus device firmware; \emph{nothing at level 2} \\
Environment ran a specific image & Attestation against the log of \S\ref{sec:log} & H
& Vendor key or firmware; meaningless without \req{M1}--\req{M3} \\
Report is not a replay & $n_v$ inside \req{B1}; ephemeral keys & C & Nothing, if
\req{V2} is performed; everything, if it is not \\
Weights released only to a measured environment & \req{K1}--\req{K2} & H & The
attestation base, plus KMS policy \\
Model commitment names the advertised weights & Merkle root over loaded bytes & C &
Second preimage on keccak \\
\textbf{The committed weights were the operands} & Receipt signed inside the
environment & \textbf{H} & Any assumption in \S\ref{sec:residual}; $sk_r$ extraction
is undetectable \\
The same, integer path & Freivalds over the opened layer (HIP-0901) & C & Nothing;
soundness is information-theoretic \\
Substitution is unprofitable & Sampled audit, bond, forfeit & E & Distinguishable
audits; bond below the run before detection \\
\textbf{Execution integrity, floating-point path} & --- & \textbf{O} & --- \\
Metadata privacy: timing, lengths, rates & --- & O & --- \\
Availability and consistency & HIP-0095 challenges & E & --- \\
\bottomrule
\end{tabular}
\caption{The classes. \req{C1} makes this table the single source every surface
projects from.}
\label{tab:class}
\end{table}

\begin{remark}[The row that matters]
The row most deployments would label ``verified inference'' is class H. It is the
claim that the committed weights were actually multiplied, and a commitment does not
establish it: a commitment binds a name to bytes, and says nothing about which bytes
were the operands of a multiplication. Where the numeric tier permits, the row below
it replaces H with C, and that is the only place in this system where the
replacement is available.
\end{remark}

% =====================================================================
\section{What a malicious operator retains}
\label{sec:residual}

Assume every requirement above is met. The operator can still do all of the
following, and a conformant implementation is one that has decided what to do about
each rather than one that has eliminated them.

\begin{description}[leftmargin=1.4em]
\item[Refuse, delay and select.] Nothing here compels service. Availability is
HIP-0095's subject; the point is that confidentiality and integrity mechanisms do
not touch it.
\item[Read the metadata.] Lengths, timings, arrival rates, client identity and
inter-token intervals are not sealed. Inter-token timing is a side channel of
practical significance for a streaming interface. Padding and batching reduce it at
a latency cost; this document requires neither and claims no metadata privacy.
\item[Attack the environment from below.] The operator controls the hypervisor, the
interrupt controller, the page tables outside the protected set, and the power
supply. The published literature against these platforms includes ciphertext side
channels from deterministic memory encryption \cite{cipherleaks}, single-stepping
frameworks \cite{sevstep}, malicious interrupt injection \cite{heckler,wesee}, and
cache-writeback manipulation \cite{cachewarp}. Each has a firmware mitigation; the
class does not close.
\item[Attack the machine physically.] The standard threat model for confidential
computing assumes a privileged software attacker and largely defers physical
attacks, on the reasonable ground that a data centre restricts physical access. Here
the adversary owns the machine, has unlimited time with it, and is paid by the
system's own reward schedule. Fault injection against the secure processor
\cite{buhren2021} and memory-aliasing attacks defeating replay protection with
inexpensive hardware \cite{badram} are published examples. Both were addressed by
platform updates; neither is the last of its kind, because the asymmetry that
produced them does not change.
\item[Extract $sk_r$, once.] This is the highest-value target in the design. One
successful extraction lets the operator author valid receipts for computations it
never performed, indefinitely, with no observable difference from honest operation.
Attestation cannot detect it: the attestation was genuine when it was made.
\req{V13} bounds the window, the distinct-measurement quorum raises the cost, and
only the audit of Section~\ref{sec:audit} detects the resulting behaviour.
\end{description}

% =====================================================================
\section{The audit, which is the only detector}
\label{sec:audit}

Every mechanism above assumes the environment behaves. The audit is the one that
does not, and it is therefore the one that detects a failure of the others.

\begin{description}[leftmargin=1.6em]
\item[\req{A1}] A conformant platform MUST audit a sampled fraction of served
requests against the reference weights it holds, on hardware the operator does not
control.

\item[\req{A2}] Audit requests MUST be indistinguishable from real ones: drawn from
the same distribution, issued through the same client identities and routing, and
opaque to the operator. A detector the operator can recognise has no power, because
the rational response is to serve honestly exactly when observed. The
confidentiality of this tier is what makes the audit undetectable, and the audit is
what detects a failure of confidentiality's own guarantees; neither is redundant.

\item[\req{A3}] Audit requests MUST be synthetic by default. Auditing a real user's
request discloses that request to the auditor, and requires the user's consent or an
auditor that is itself a measured environment.

\item[\req{A4}] Where $\kappa$ declares batch-invariant kernels and the seed is
committed per \req{R2}, the audit MUST be an exact replay and an equality check.
Where it is not, the audit is a statistical test with a false-positive rate and a
power curve, and every surface reporting its result MUST describe it as a detector
rather than as a proof.

\item[\req{A5}] The audit fraction $\alpha$, the per-audit detection probability
$\beta$ and the bond $S$ MUST satisfy $\alpha \beta S > \Delta$, where $\Delta$ is
the cost an operator saves per substituted request. Because $\Delta$ is a marginal
serving cost and $S$ is a bond, $\alpha$ can be small. The assumptions carry more
weight than the inequality: the operator is risk-neutral, cannot distinguish audited
requests, forfeits with certainty on detection, and has not already withdrawn the
bond. The inequality also bounds expected value and not damage --- the number of
substitutions before the first detection is geometric with mean
$(\alpha\beta)^{-1}$ --- so where harm is not linear in $\Delta$, $S$ MUST cover
that run and not one request. HIP-0095 and HIP-0096 set these numbers; this document
states the constraint they satisfy.
\end{description}

% =====================================================================
\section{Conformance}
\label{sec:conformance}

An implementation claims conformance for one of three roles, and the roles are
disjoint.

\begin{description}[leftmargin=1.4em]
\item[Environment.] The measured image. Conforms if it satisfies \req{B1},
\req{B1a}, \req{B2}--\req{B5}, \req{B7}--\req{B8}, \req{E1}--\req{E3}, \req{K3} and
\req{R1}--\req{R3}.
\item[Verifier.] Whatever decides a bundle is acceptable --- a key broker, a
scheduler, a chain-side voucher, a user's client. Conforms if it satisfies
\req{V0}--\req{V15} and \req{L2}. A verifier is the only role permitted to assign a
level.
\item[Settlement.] Whatever pays. Conforms if it satisfies \req{S1}--\req{S5} and
consumes only verifier-assigned levels.
\end{description}

\begin{description}[leftmargin=1.6em]
\item[\req{X1}] A deployment MUST refuse to settle at this tier until a verifier
conforming to \req{V0}--\req{V15} exists for the platform in question.
\emph{Observed:} the deployed configuration already refuses --- \code{CCEvidence}'s
voucher set is empty and \code{TEEEvidence}'s threshold defaults to zero, which the
contract treats as never attesting. Nothing mis-settles today; the work is
unfinished rather than broken, and \req{X1} says the order in which it may be
finished.

\item[\req{X2}] There MUST be one attester and one verifier. \emph{Observed:} four
places hold something that answers to the name today, and three of them answer
\emph{yes} to everything. \code{net/hanzo-pqc}, \code{rust-sdk/crates/hanzo-kbs} and
\code{rust-sdk/crates/hanzo-crypto} each carry a \code{MockAttestationVerifier}:
two return success for any quote against any expected measurement, the third
accepts any non-empty report as an SEV-SNP attestation at level 2, and two of them
additionally return a fabricated H100 result with confidential mode reported on. Each
enumerates the real steps in comments. The fourth, \code{hanzo-attest} in
\code{hanzoai/node}, declares the two hardware roots and refuses them by name.

\emph{Proposed:} the implementation belongs where the refusal already is, because
refusing by name is the correct behaviour for an unimplemented root and returning
success is not. The three mocks are deleted at the same time rather than deprecated.
A mock that returns success is not a placeholder for a verifier; it is a verifier
that accepts every forgery, and it is reachable from anything that takes the trait.

\item[\req{X3}] The vendor root pool MUST hold the root of every vendor whose
reports the deployment accepts, MUST be embedded rather than read from a host trust
store, and MUST NOT hold a root for a platform the verifier cannot fully check.
\emph{Observed:} exactly one root is embedded, Intel's SGX provisioning root, so
AMD SEV-SNP and NVIDIA device chains fail closed. That is correct behaviour and an
incomplete pool: it means no SEV-SNP host and no confidential accelerator can attest
at all today.
\end{description}

% =====================================================================
\section{Unsolved}
\label{sec:open}

Stated as problems, not as a plan. The companion paper \cite{paper} argues each; this
document records that no requirement above closes them.

\begin{description}[leftmargin=1.4em]
\item[Execution integrity in floating point.] The gap between the integer path,
where Freivalds gives an information-theoretic guarantee, and the floating-point
path, where nothing does, is the central open problem of this tier. A tolerance wide
enough to admit honest reductions is wide enough to hide a substitution, which is
why HIP-0901 routes by numeric tier rather than weakening the check, and why this
tier exists at all.
\item[Detecting key extraction.] $sk_r$ compromise is undetectable by construction.
\req{V13} bounds the window and \req{A1} detects the behaviour; neither detects the
extraction.
\item[Physical attack by the party being attested.] Whether any composition of
attestation, tamper evidence and bonding meaningfully raises the cost, or whether the
honest answer is that certain models do not go on third-party hardware, is a policy
question \req{K4} defers to a human rather than resolving.
\item[Metadata.] Nothing here protects timing, lengths or rates.
\item[Post-quantum attestation.] No post-quantum device-identity root exists in
silicon. \req{C4} requires saying so.
\end{description}

% =====================================================================
\appendix
\section{The tree against this document}
\label{app:tree}

Observed at the time of writing, so that conformance is measured rather than
asserted. \emph{Present} means the code exists and does what this document requires
of it; \emph{partial} means it exists and diverges; \emph{absent} means it does not
exist.

One row moved while this was being written, and it is worth recording because it is
the shape the rest should take. \code{hanzo-attest} previously folded its own
measurement over the boot inputs, which would have made two answers to the question
of what a launch is --- one in the node, one in the monitor, drifting apart the first
time the monitor changed what it hands the hypervisor. It now reads the monitor's
document and signs the 64 bytes that document names. The register, the extension rule
and the digest are the monitor's, which is what \req{B7} pins them to.

\begin{center}
\small
\begin{tabular}{@{}L{0.27\textwidth}L{0.25\textwidth}lL{0.27\textwidth}@{}}
\toprule
Requirement & Where & State & Divergence \\
\midrule
\req{B7}, \req{B8} registers & \code{hanzoai/vm}, \code{vm-measure} & present & --- \\
\req{B1} binding & same & partial & Binds $L \,\|\, W$ only; no keys, no nonce \\
One measurement, \req{B7} & \code{hanzoai/node}, \code{hanzo-attest} & present &
Reads the monitor's document and signs the 64 bytes it names; folds no second
register of its own \\
\req{B6} digest choice & same & present & --- \\
\req{E2} raw report & same, \code{attest.rs} & present & Request ABIs and ioctl
encodings implemented and pinned by test; never executed, for want of hardware \\
\req{E3} no synthesis & same & present & --- \\
\req{E3} no synthesis & \code{hanzoai/node}, \code{tee\_attestation.rs} &
partial & Simulation and development modes fabricate reports; production falls back
to simulation \\
\req{V0}--\req{V15} & \code{hanzo-pqc}, \code{hanzo-kbs}, \code{hanzo-crypto} &
absent & Three mock verifiers return success for any input \\
\req{V0}--\req{V15} & \code{luxfi/precompile} \code{0x0301} & partial & Validates a
chain to an embedded root over a 48-byte header carrying no measurement, policy, TCB
or report-data field \\
\req{L1} one enumeration & \code{hanzo-pqc} & partial & \code{PrivacyTier} and
\code{CCTier} both present \\
\req{L2} derived level & same & partial & \code{max\_tier} reads self-reported
features \\
\req{M1}--\req{M4} log & --- & absent & No build-to-measurement mapping is published \\
\req{R1}--\req{R3} receipt & \code{hanzo-engine} transcript & absent & Transcript
machinery present; no receipt emitted from a live forward pass \\
\req{S1}--\req{S5} & \code{luxfi/standard} \code{contracts/ai/compute/} & partial &
Backends present and fail closed; voucher set empty, threshold zero \\
\req{A1}--\req{A5} audit & --- & absent & No canary traffic, no reference check, no
batch-invariant audit path \\
Confidential launch path & \code{hanzoai/vm} & absent & The microVM layer boots k3s
on an ordinary hypervisor; a confidential guest is a different launch path with
different firmware and a measurement to reproduce \\
\bottomrule
\end{tabular}
\end{center}

Two observations about this table rather than about its rows.

The audit is absent, and it is the only mechanism here that detects a compromised
environment rather than assuming one cannot occur. That argues for building it before
the parts that assume it away, which is the reverse of the order the table's other
absences suggest.

No requirement in this document has been exercised on confidential hardware. Both
hosts on the bench were checked directly at the time of writing. One is x86-64, an
AMD Ryzen AI Max+ 395, whose \code{/proc/cpuinfo} carries no \code{sev},
\code{sev\_es}, \code{sev\_snp} or \code{tdx} flag and which has no \code{/dev/sev}
device. The other is aarch64 with an NVIDIA GB10 and no guest device of either kind.
Neither absence is a configuration that can be changed: one is client silicon, the
other an Arm part, and confidential guests on both platforms require server parts
with the feature fused on. Exercising this tier therefore requires hardware the bench
does not have. It is specified and unexercised, and no performance claim is made for
it.

% =====================================================================
\begin{thebibliography}{99}

\bibitem{paper}
Hanzo AI Research.
\newblock Attested inference on operator hardware: what confidential computing
establishes for a third-party AI operator, what the settlement commitment adds, and
what remains unsolved.
\newblock Hanzo Papers, \code{papers/attestation}.

\bibitem{hanzod}
Hanzo AI Research.
\newblock Attested compute in \code{hanzod}: boot measurement, host-asserted
evidence, and the gap to hardware attestation.
\newblock Hanzo Papers.

\bibitem{hipsettle}
Hanzo AI.
\newblock Compute settlement: credits in, attested work out, on a Hanzo sovereign L1
over Lux.
\newblock Hanzo Improvement Proposal, number assigned on merge.

\bibitem{hip0901}
Hanzo AI.
\newblock HIP-0901: Proof of AI --- native execution proofs, canonical contract, and
operator-LLM governance.

\bibitem{hip0095}
Hanzo AI.
\newblock HIP-0095: QoS challenge system.

\bibitem{hip0096}
Hanzo AI.
\newblock HIP-0096: AI compute contribution rewards.

\bibitem{hip0027}
Hanzo AI.
\newblock HIP-0027: Secrets management standard.

\bibitem{hip1325}
Hanzo AI.
\newblock HIP-1325: Node --- your machines on a socket.

\bibitem{snpabi}
Advanced Micro Devices.
\newblock \emph{SEV Secure Nested Paging Firmware ABI Specification}.

\bibitem{tdx}
Intel Corporation.
\newblock \emph{Intel Trust Domain Extensions (TDX) Module Architecture
Specification} and \emph{Data Center Attestation Primitives}.

\bibitem{hpke}
R. Barnes, K. Bhargavan, B. Lipp, and C. Wood.
\newblock Hybrid Public Key Encryption.
\newblock RFC 9180, 2022.

\bibitem{rfc6962}
B. Laurie, A. Langley, and E. Kasper.
\newblock Certificate Transparency.
\newblock RFC 6962, 2013.

\bibitem{buhren2021}
R. Buhren, H.-N. Jacob, T. Krachenfels, and J.-P. Seifert.
\newblock One glitch to rule them all: fault injection attacks against AMD's Secure
Encrypted Virtualization.
\newblock \emph{ACM Conference on Computer and Communications Security}, 2021.

\bibitem{cipherleaks}
M. Li, Y. Zhang, H. Wang, K. Li, and Y. Cheng.
\newblock CipherLeaks: breaking constant-time cryptography on AMD SEV via the
ciphertext side channel.
\newblock \emph{USENIX Security Symposium}, 2021.

\bibitem{sevstep}
L. Wilke, J. Wichelmann, A. Rabich, and T. Eisenbarth.
\newblock SEV-Step: a single-stepping framework for AMD-SEV.
\newblock 2024.

\bibitem{heckler}
B. Schl\"uter, S. Sridhara, A. Bertschi, and S. Shinde.
\newblock Heckler: breaking confidential VMs with malicious interrupts.
\newblock 2024.

\bibitem{wesee}
B. Schl\"uter, S. Sridhara, M. Kuhne, A. Bertschi, and S. Shinde.
\newblock WeSee: using malicious \#VC interrupts to break AMD SEV-SNP.
\newblock 2024.

\bibitem{cachewarp}
R. Zhang, L. Gerlach, D. Weber, L. Hetterich, Y. Lu, A. Kogler, and M. Schwarz.
\newblock CacheWarp: software-based fault injection using selective state reset.
\newblock \emph{USENIX Security Symposium}, 2024.

\bibitem{badram}
J. De Meulemeester, L. Wilke, D. Oswald, T. Eisenbarth, I. Verbauwhede, and
J. Van Bulck.
\newblock BadRAM: practical memory aliasing attacks on trusted execution
environments.
\newblock 2025.

\end{thebibliography}

\end{document}
